\documentclass[11pt]{amsart}

\usepackage[T1]{fontenc}
\usepackage{lmodern}
\usepackage{microtype}
\usepackage{amsmath,amssymb,amsthm}
\usepackage[hidelinks]{hyperref}

\newtheorem{theorem}{Theorem}

\newcommand{\Av}{\operatorname{Av}}
\newcommand{\inv}{\operatorname{inv}}

\title[Frankl\'{\i}n's conjecture on $(321,1342)$-avoiders]
{A Counterexample to a Conjecture of Frankl\'{\i}n on
$(321,1342)$-Avoiders}

\subjclass[2020]{05A05, 05A15}
\keywords{permutation patterns, inversions, indecomposable permutations}
\date{}

\begin{document}

\begin{abstract}
Frankl\'{\i}n conjectured that the number of indecomposable permutations
avoiding $321$ and $1342$ with exactly $k$ inversions is
$1+\binom{k+1}{2}$. The case $k=1$ already gives a counterexample.
Using a disjoint structural cover of $\Av(321,1342)$, we determine the
bivariate generating function of the indecomposable members by length
and inversions and prove that the exact count is $1+\binom{k}{2}$, so
the conjectured sequence is the correct one shifted by one in the index.
The number of indecomposable avoiders of length $n$ is $1+\binom{n}{3}$.
\end{abstract}

\maketitle

\section{Exact enumeration}

For a set $B$ of patterns, let $\Av(B)$ denote the permutations avoiding
every pattern in $B$, and let $\inv(\pi)$ be the inversion number of
$\pi$. A nonempty permutation $\pi=\pi_1\cdots\pi_n$ is
\emph{indecomposable} if no proper prefix $\pi_1\cdots\pi_j$ has value
set $\{1,\ldots,j\}$. Define
\[
I_k=
\{\pi\in\Av(321,1342):
  \pi\text{ is indecomposable and }\inv(\pi)=k\},
\]
and let $S_n^*$ be the indecomposable members of $\Av(321,1342)$ of
length $n$. These are the sets written $I_k(321,1342)$ and
$S_n^*(321,1342)$ in \cite{Franklin2025}, where the bare symbols $I_k$
and $S_n^*$ denote the corresponding sets of all indecomposable
permutations; we use the abbreviated symbols throughout. As recorded in
\cite{Franklin2025}, an indecomposable permutation of length $n$ has at
least $n-1$ inversions, so $I_k$ consists of permutations of length at
most $k+1$ and is in particular finite; this is what makes enumeration
by inversions possible at all. For the members of $\Av(321,1342)$ the
same bound is recovered in the proof of Theorem~\ref{thm:main}.

In an unnumbered remark in the introduction of
\cite[p.~3]{Franklin2025}, offered on the basis of numerical evidence,
Frankl\'{\i}n writes: ``Another case is $I_k(321,1342)$, which we
conjecture to have $k(k+1)/2+1$ elements.'' In the abbreviated notation
this is the assertion
\[
|I_k|=1+\binom{k+1}{2}.
\]
By Theorem~\ref{thm:main}, the conjectured value $1+\binom{k+1}{2}$ is
precisely $|I_{k+1}|$, so the conjectured sequence is the correct one
shifted by one in the index. The shift is not a difference of
convention: in \cite{Franklin2025} as here, $k$ is the number of
inversions of the permutation and $I_0=\{1\}$.

A permutation with one inversion is obtained from the identity by
interchanging two adjacent consecutive values, and it is indecomposable
only in the case $21$. Hence $I_1=\{21\}$, whereas the conjectured
formula gives $|I_1|=2$.

The same passage of \cite{Franklin2025} reports experimental evidence
that $S_n^*(321,1342)$ grows like $n^3/6$; the second formula of
Theorem~\ref{thm:main} confirms that prediction and sharpens it to the
closed form $1+\binom{n}{3}$. The asymptotic bound on $S_n^*(321,1342)$
drawn there from the conjectured count is unaffected by the correction,
since $1+\binom{k}{2}$ and $1+\binom{k+1}{2}$ are both asymptotic to
$k^2/2$. The following theorem gives the exact correction.

\begin{theorem}\label{thm:main}
Let
\[
J(x,q)=
\sum_{\substack{\pi\in\Av(321,1342)\\
                 \pi\text{ indecomposable}}}
x^{|\pi|}q^{\inv(\pi)}.
\]
Then
\begin{equation}\label{eq:J}
J(x,q)
=
\frac{x}{1-xq}
+
\frac{x^3q^2}{1-xq}
\sum_{a\ge0}
\frac{(xq)^a}
     {(1-xq^{a+1})(1-xq^{a+2})}.
\end{equation}
Consequently, for every $k\ge0$ and $n\ge1$,
\[
|I_k|=1+\binom{k}{2},
\qquad
|S_n^*|=1+\binom{n}{3}.
\]
In particular, the conjectured formula is false for every $k\ge1$.
\end{theorem}

\begin{proof}
The disjoint structural cover of Bean, Gudmundsson, and Ulfarsson
\cite[Proposition~7 and Equation~(6)]{BeanGudmundssonUlfarsson2019}
shows that every nonempty member of $\Av(321,1342)$ has a unique
representation of one of the following two types. The cover is stated
there as a gridded tiling, in which the maximum shares a column with an
increasing cell and hence interleaves it; we have unpacked it into the
explicit block form below. All displayed blocks may be empty, and
inequalities involving blocks are understood elementwise.

In the first type,
\[
\pi=\alpha\,n\,\beta,
\]
where $\alpha\in\Av(321,1342)$, the block $\beta$ is increasing, and
all entries of $\alpha$ are smaller than all entries of $\beta$. If
$\alpha$ is nonempty, its entries are $1,\ldots,|\alpha|$, so $\pi$ is
decomposable. Thus the indecomposable permutations of this type are
\[
n\,1\,2\cdots(n-1),\qquad n\ge1,
\]
and their contribution to $J(x,q)$ is
\[
\sum_{n\ge1}x^nq^{n-1}
=
\frac{x}{1-xq}.
\]

In the second type,
\[
\pi=c\,\gamma_1\gamma_2\,n\,b\,\delta_1\delta_2,
\]
where the four blocks are increasing and
\[
\gamma_2<b<\delta_1<c<\gamma_1<\delta_2<n.
\]
Every such permutation is indecomposable. Indeed, a proper prefix
ending before $n$ contains $c$ but omits the smaller entry $b$, while a
proper prefix of length $j<n$ containing $n$ contains a value larger
than $j$.

Set
\[
a=|\gamma_1|,
\qquad
u=|\gamma_2|,
\qquad
v=|\delta_1|,
\qquad
w=|\delta_2|.
\]
The length is $3+a+u+v+w$. The inversions are those from $c$ to
$\gamma_2,b,\delta_1$, from each entry of $\gamma_1$ to
$\gamma_2,b,\delta_1$, and from $n$ to $b,\delta_1,\delta_2$. Hence
\[
\inv(\pi)
=
(u+1+v)+a(u+1+v)+(1+v+w)
=
2+a(u+v)+a+u+2v+w.
\]
Therefore the second type contributes
\begin{align*}
&\sum_{a,u,v,w\ge0}
x^{3+a+u+v+w}
q^{2+a(u+v)+a+u+2v+w}
\\
&\qquad=
\frac{x^3q^2}{1-xq}
\sum_{a\ge0}
\frac{(xq)^a}
     {(1-xq^{a+1})(1-xq^{a+2})},
\end{align*}
which proves \eqref{eq:J}.

As a check on the transcription, the two types give, ignoring
inversions, the functional equation
\[
F(x)
=
1
+
\frac{xF(x)}{1-x}
+
\frac{x^3}{(1-x)^4}
\]
for the generating function $F$ of $\Av(321,1342)$ by length, the
summand $1$ accounting for the empty permutation. Hence
\begin{align*}
F(x)
&=
\frac{(1-x)^4+x^3}{(1-x)^3(1-2x)}
=
\frac{1-4x+6x^2-3x^3+x^4}{1-5x+9x^2-7x^3+2x^4}
\\
&=
1+x+2x^2+5x^3+13x^4+32x^5+\cdots,
\end{align*}
which is the generating function of the class recorded in
\cite{BeanGudmundssonUlfarsson2019}.

For the first type, $\inv(\pi)=|\pi|-1$; for the second,
\[
\inv(\pi)-(|\pi|-1)
=
a(u+v)+v
\ge0.
\]
Thus setting $x=1$ is well defined coefficientwise. For every $A\ge0$,
\[
\sum_{a=0}^{A}
\frac{q^a}{(1-q^{a+1})(1-q^{a+2})}
=
\frac{1}{1-q}
\left(
\frac{1}{1-q}
-
\frac{q^{A+1}}{1-q^{A+2}}
\right).
\]
Letting $A\to\infty$ coefficientwise gives
\[
\sum_{a\ge0}
\frac{q^a}{(1-q^{a+1})(1-q^{a+2})}
=
\frac{1}{(1-q)^2}.
\]
Consequently,
\[
J(1,q)
=
\frac{1}{1-q}
+
\frac{q^2}{(1-q)^3},
\]
and coefficient extraction yields
\[
|I_k|
=
1+\binom{k}{2}.
\]
Finally, setting $q=1$ in \eqref{eq:J} gives
\[
J(x,1)
=
\frac{x}{1-x}
+
\frac{x^3}{(1-x)^4},
\]
whose coefficient of $x^n$ is $1+\binom{n}{3}$.
\end{proof}

\begin{thebibliography}{9}

\bibitem{BeanGudmundssonUlfarsson2019}
C.~Bean, B.~A.~Gudmundsson, and H.~\'{U}lfarsson,
\emph{Automatic discovery of structural rules of permutation classes},
Math. Comp. \textbf{88} (2019), no.~318, 1967--1990.
\href{https://doi.org/10.1090/mcom/3386}
{doi:10.1090/mcom/3386}.

\bibitem{Franklin2025}
A.~F.~Frankl\'{\i}n,
\emph{Pattern avoiding permutations enumerated by inversions},
Discrete Math. Theor. Comput. Sci. \textbf{27} (2025), no.~1,
Paper No.~5, 21~pp.
\href{https://doi.org/10.46298/dmtcs.14437}
{doi:10.46298/dmtcs.14437}.

\end{thebibliography}

\end{document}